\chapter{序論}\label{chap:intro}

\section{研究背景}\label{sec:background}

近年，〇〇の研究が盛んである\cite{einstein1905,knuth1984}．
特に\textbf{太字}や\textit{イタリック}，\emph{強調}，\texttt{等幅}も使える．
脚注も書ける\footnote{これは脚注です．}．

\subsection{箇条書き}

\begin{itemize}
  \item 一つ目の項目
  \item 二つ目の項目
  \begin{enumerate}
    \item 入れ子の番号付き
    \item もう一つ
  \end{enumerate}
\end{itemize}

\section{本論文の構成}

本論文の構成は以下のとおりである．
第\ref{chap:method}章の式\eqref{eq:euler}と図\ref{fig:sample}，表\ref{tab:result}を参照．
節\ref{sec:background}も参照できる．
